\documentclass[a4paper,11pt]{article}
\usepackage[utf8]{inputenc}
\usepackage[T1]{fontenc}
\usepackage[french]{babel}
\usepackage[left=1cm,right=1cm,top=.5cm,bottom=0cm]{geometry}
\usepackage{enumitem}
\usepackage{hyperref}
\usepackage{xcolor}
\usepackage{titlesec}
\usepackage{fontawesome5}
\usepackage{lmodern}
\usepackage[normalem]{ulem}

% --- Couleurs ---
\definecolor{primary}{RGB}{35, 55, 123}
\definecolor{secondary}{RGB}{80, 80, 80}

% --- Configuration des liens ---
\hypersetup{colorlinks=true, linkcolor=primary, filecolor=primary, urlcolor=primary}

% --- Formatage des titres ---
\titleformat{\section}{\large\bfseries\color{primary}\uppercase}{}{0em}{}[\titlerule]
\titlespacing{\section}{0pt}{8pt}{5pt}

% --- Commandes personnalisées ---
\newcommand{\resumeItem}[1]{\item\small{#1}}
\newcommand{\resumeSubheading}[4]{
  \vspace{-1pt}\item
    \begin{tabular*}{0.97\textwidth}[t]{l@{\extracolsep{\fill}}r}
      \textbf{#1} & \textit{\small #2} \\
      \textit{\small #3} & \textit{\small #4} \\
    \end{tabular*}\vspace{-5pt}
}

\begin{document}

% --- EN-TÊTE ---
\begin{center}
    {\Large \textbf{Loïc Aron MBASSI EWOLO}} \\ \vspace{4pt}
    \textbf{\large Stage Développement \& Analyse de Performance - Calcul Scientifique} \\
    \textit{\small Disponible dès \textbf{Avril 2026} pour 4 à 6 mois} \\ \vspace{4pt}
    \small
    \faMapMarker\ Cergy, France (Mobile) \quad
    \faPhone\ (+33) 7 59 52 33 98 \quad
    \faEnvelope\ \href{mailto:loicaron.mbassiewolo@ensea.fr}{loicaron.mbassiewolo@ensea.fr} \\ \vspace{2pt}
    \faLinkedin\ \href{https://www.linkedin.com/in/loic-aron-mbassi-ewolo-3942a9246/}{linkedin} \quad
    \faGithub\ \href{https://github.com/Nameless0l}{github} \quad
    \faGlobe\ \href{https://mbassiloic.tech}{mbassiloic.tech}
\end{center}

\vspace{-6pt}

% --- PROFIL ---
\noindent\small{Élève-ingénieur en double diplôme \textbf{ENSEA/Polytechnique Yaoundé}, avec une expérience en programmation \textbf{C/C++}, \textbf{Python} et environnement \textbf{Linux}. Projets en programmation système (IPC, mémoire partagée) et optimisation sur cibles embarquées. Je recherche un stage au \textbf{CEA Grenoble} pour contribuer au développement de micro-programmes pour le calcul scientifique.}

% --- FORMATION ---
\section{Formation}
\begin{itemize}[leftmargin=*, label=\tiny$\bullet$, nosep]
    \item \textbf{ENSEA (École Nationale Supérieure de l'Électronique et de ses Applications)} \hfill \textit{Cergy, France | 2025 -- Présent} \\
    \small{Ingénieur 2A, Double Diplôme - Systèmes Embarqués, Traitement du Signal et IA.} \\
    \textit{\small{Cours : Programmation C/C++, Architecture des processeurs, Systèmes temps réel, Linux embarqué.}}
    \vspace{2pt}
    \item \textbf{École Polytechnique de Yaoundé (1\textsuperscript{ère} école d'ingénieur CEMAC)} \hfill \textit{Cameroun | 2021 -- 2025} \\
    \small{Diplôme d'Ingénieur de Conception (Master 2) : \textbf{Mathématiques Appliquées}, Algorithmique, Génie Logiciel.}
    \item \textbf{Université de Yaoundé I} \hfill \textit{Cameroun | 2020 -- 2022} \\
    \small{Licence Informatique \& Mathématiques : Algorithmique C, Analyse Numérique, Algèbre Linéaire.}
\end{itemize}

% --- COMPÉTENCES TECHNIQUES ---
\section{Compétences Techniques}
\begin{itemize}[leftmargin=*, nosep]
    \item \small{\textbf{Langages :} \textbf{C/C++}, \textbf{Python}, Bash/Shell scripts, Java.}
    \item \small{\textbf{Systèmes :} \textbf{Linux} (Ubuntu, Debian), Programmation système (IPC, SHM, Signaux), Docker.}
    \item \small{\textbf{Calcul \& Optimisation :} NumPy, PyTorch, Optimisation sur GPU (Jetson), Algorithmes numériques.}
    \item \small{\textbf{Mathématiques :} Algèbre Linéaire, Analyse Numérique, Probabilités/Statistiques, Optimisation.}
    \item \small{\textbf{Outils :} Git, CMake, GDB, Valgrind, CI/CD.}
\end{itemize}

% --- PROJETS SYSTÈMES ---
\section{Projets en Programmation Système \& Optimisation}
\begin{itemize}[leftmargin=*, label={-}]

\item \textbf{Projet TAXI - Simulation Multi-Processus en C} \hfill \textit{Projet Académique | 2024}
\vspace{-2pt}
\begin{itemize}[leftmargin=0.4cm, nosep, label=\tiny$\bullet$]
    \item \small{Développement d'un noyau de simulation pour une flotte de taxis sous \textbf{Linux}.}
    \item \small{Gestion de la \textbf{mémoire partagée} (SHM), communication inter-processus (IPC), synchronisation par signaux.}
    \item \small{Implémentation d'un ordonnanceur FIFO avec prévention des interblocages.}
    \item \small{Visualisation temps réel avec OpenGL.}
\end{itemize}
\vspace{3pt}

\item \textbf{Upgrade Jetson - Optimisation IA Embarquée (Défense)} \hfill \textit{Challenge CoHoMa | En cours}
\vspace{-2pt}
\begin{itemize}[leftmargin=0.4cm, nosep, label=\tiny$\bullet$]
    \item \small{Optimisation d'algorithmes de SLAM et de détection (YOLOv10) sur \textbf{Nvidia Jetson}.}
    \item \small{Analyse de performance et profilage pour respecter les contraintes temps réel.}
    \item \small{Utilisation de \textbf{Docker} pour l'isolation et la reproductibilité des environnements.}
\end{itemize}
\vspace{3pt}

\item \textbf{Stage Recherche IoT Lab - TinyML} \hfill \textit{ENSPY | Oct. 2023 - Jan. 2024}
\vspace{-2pt}
\begin{itemize}[leftmargin=0.4cm, nosep, label=\tiny$\bullet$]
    \item \small{Optimisation de modèles ML pour \textbf{hardware limité} (Arduino, Raspberry Pi).}
    \item \small{Travail sous contraintes de mémoire et de temps d'exécution.}
    \item \small{Technologies : Python, TensorFlow Lite, C.}
\end{itemize}
\vspace{3pt}

\item \textbf{GoFind - Architecture Microservices} \hfill \textit{Projet Académique | 2024}
\vspace{-2pt}
\begin{itemize}[leftmargin=0.4cm, nosep, label=\tiny$\bullet$]
    \item \small{Conception d'une architecture microservices avec communication asynchrone (RabbitMQ).}
    \item \small{Déploiement avec \textbf{Docker}, scripts d'automatisation en \textbf{Bash}.}
\end{itemize}

\end{itemize}

% --- EXPÉRIENCE RECHERCHE ---
\section{Expérience en Recherche}
\begin{itemize}[leftmargin=*, label={-}]

    \resumeSubheading
      {Assistant de Recherche - Généralisation des Réseaux de Neurones}{Juin 2025 -- Sept. 2025}
      {MILA (Institut Québécois d'IA) - Remote}{Québec, Canada}
      \begin{itemize}[leftmargin=0.4cm, nosep, label=\tiny$\bullet$]
        \item \small{Implémentation de pipelines de tests sous \textbf{Python/PyTorch}.}
        \item \small{Reproduction d'expériences et analyse de convergence mathématique.}
        \item \small{Travail en autonomie avec reporting régulier.}
      \end{itemize}

\end{itemize}

% --- CERTIFICATIONS & LANGUES ---
\section{Certifications, Leadership \& Langues}
\begin{itemize}[leftmargin=*, nosep]
    \item \textbf{Certifications :} Supervised ML (DeepLearning.AI/Stanford), Python for Data Science (IBM).
    \item \textbf{Leadership :} Chef de la Cellule Projets (Polytechnique Yaoundé), animation de workshops techniques.
    \item \textbf{Hackathons :} \textbf{1\textsuperscript{ère} Place} CITS National (Algorithmes d'optimisation), IndabaX Africa.
    \item \textbf{Langues :} Français (Natif), \textbf{Anglais} (Intermédiaire/Technique).
\end{itemize}

\end{document}
